\chapter{Gyrokinetics and Multiscale Reduction}\label{ch:19}
\begin{chaptersummary}
The strong-field ordering of Chapter 4 separates fast gyromotion from slow transport. Coordinate transformations and ordered averaging remove the fast angle while retaining parallel-velocity and magnetic-moment distributions. Gyroaveraging and reconstruction of physical moments reconnect reduced particles with self-consistent fields. A uniform-field electrostatic model makes this five-dimensional reduction concrete.
\end{chaptersummary}
\learninggoals{Specify strong-field orderings and five-dimensional coordinates; distinguish guiding centers from gyrocenters; calculate Fourier gyroaverages; separate kinetic reduction, moment closure, field coupling, and numerical discretization.}
\studyroute{A controlled uniform-field electrostatic example is given; general electromagnetic Lie transforms, second-order Hamiltonians, magnetization, and collision transformations in magnetic geometry are omitted. A complete model must use a consistent variational order.}{guiding center; gyrocenter; gyroaverage; finite Larmor radius; gyrokinetic polarization; electromagnetic gyrokinetics}{Gyroaverage a perpendicular Fourier potential in a uniform field, derive the Bessel factor, and write linear quasineutrality for a specified nonadiabatic distribution variable.}

\section{Scale Ordering and the Five-Dimensional State}\index{Gyrokinetics}
Using equilibrium scale $L_B$ and ion thermal gyroradius $\rho_{Li}$, a common perturbation ordering is
\[
\epsilon=\rho_{Li}/L_B\ll1,\qquad
\omega/\Omega_i=O(\epsilon),\qquad
k_\parallel/k_\perp=O(\epsilon),\qquad
k_\perp\rho_{Li}=O(1).
\]
Weak perturbations may additionally satisfy $|q_i\delta\phi|/T_i=O(\epsilon)$.
Below, $T_s$ denotes temperature in energy units, corresponding to $k_BT_s$ earlier.
Electromagnetic, collisional, and electron-scale models require additional orderings.

The original phase space $(\bm x,\bm v)$ can be transformed to
$(\bm R,v_\parallel,\mu,\vartheta)$,
where $\mu=m v_\perp^2/(2B_0)$ and $\vartheta$ is the fast gyrophase.
A guiding-center transformation first handles rapid gyromotion in the equilibrium field;
a gyrocenter transformation then handles perturbation-induced angle dependence.
The reduced distribution $F_s(\bm R,v_\parallel,\mu,t)$ no longer depends independently on $\vartheta$,
retaining three position and two velocity coordinates.
For uniform $B_0$ and a gyrophase-independent function $A$, velocity moments use
\[
\int A\dd^3v=\frac{2\pi B_0}{m_s}
\int_{-\infty}^{\infty}\int_0^\infty A\dd\mu\dd v_\parallel.
\]
The angle integral has been performed analytically; parallel velocity and magnetic moment remain independent variables.

\conceptbridge{Why do field equations still require polarization and gyroaveraging after a fast angle has been removed?}
{The position coordinate changes from particle position to gyrocenter, but density and current are measured at physical positions.
Pushing the reduced distribution back to physical space produces finite-Larmor-radius, polarization, and magnetization contributions.
These arise at the same transformation order as the reduced dynamics and determine consistent fields and energy pairings.}
{Phase-space transformation $\longrightarrow$ five-dimensional evolution $\longrightarrow$ physical moments and self-consistent fields.}

\section{Gyroaveraging and Finite Larmor Radius}
Write particle position as $\bm x=\bm R+\bm\rho(\vartheta)$.
At fixed gyrocenter and magnetic moment, define
\[
\langle\phi\rangle_{\bm R}
=\frac1{2\pi}\int_0^{2\pi}
\phi(\bm R+\bm\rho(\vartheta))\dd\vartheta.
\]
For a plane wave in a uniform magnetic field,
\[
\left\langle\ee^{\ii\bm k\cdot\bm x}\right\rangle_{\bm R}
=\ee^{\ii\bm k\cdot\bm R}
J_0(k_\perp v_\perp/\Omega_s).
\]
finite $k_\perp\rho_{Ls}$ changes the particle response through ring averaging.
The small parameter removes the fast angle, while the Bessel factor retains its influence on slow-scale coupling.

\section{Kinetic Structure and a Specific Model}
General gyrocenter dynamics can be written in characteristic form:
\[
\partial_tF_s+\dot{\bm R}_s\cdot\nabla_{\bm R}F_s
+\dot v_{\parallel s}\partial_{v_\parallel}F_s=C_s^{\mathrm{GK}}[F].
\]
Characteristics, phase-space measure, and collisions must use the same reduced coordinates;
a nonuniform field requires the corresponding phase-space Jacobian in conservative form.

As a computable example, take uniform $B_0\bm b$, no equilibrium gradients, no collisions, and low-frequency electrostatic perturbations.
Let $F_{0s}$ be a uniform Maxwellian at temperature $T_s$,
and $h_s$ the usual nonadiabatic distribution, with its convention fixed by the physical-density formula below.
A corresponding nonlinear electrostatic system is
\begin{equation}\label{eq:slab-gk}
\partial_th_s+v_\parallel\partial_\parallel h_s
+\bm v_{E,s}\cdot\nabla_\perp h_s
=\frac{q_sF_{0s}}{T_s}\partial_t\langle\phi\rangle_{\bm R},
\qquad
\bm v_{E,s}=\frac{\bm b\times\nabla\langle\phi\rangle_{\bm R}}{B_0}.
\end{equation}
The drift is the $\bm E\times\bm B/B_0^2$ velocity of the ring-averaged electric field.
A uniform equilibrium has no gradient drive; density or temperature gradients require matching drive terms.
For general electromagnetic construction and energy conservation, see~\cite{brizard2007,abel2013}.

\section{Quasineutrality, Polarization, and Amp\`ere Coupling}
For a Fourier mode of this uniform-field model, the physical density perturbation is
\[
\delta n_{s,\bm k}
=-\frac{q_sn_{0s}}{T_s}\phi_{\bm k}
+\int J_{0s}h_{s,\bm k}\dd^3v,\qquad
J_{0s}=J_0(k_\perp v_\perp/\Omega_s).
\]
Thus quasineutrality, $\sum_sq_s\delta n_s=0$, gives
\begin{equation}\label{eq:gk-qn}
\sum_s\frac{q_s^2n_{0s}}{T_s}\phi_{\bm k}
=\sum_sq_s\int J_{0s}h_{s,\bm k}\dd^3v.
\end{equation}
If instead $g_{s,\bm k}=h_{s,\bm k}-(q_sF_{0s}/T_s)J_{0s}\phi_{\bm k}$,
the same relation becomes
\[
\sum_s\frac{q_s^2n_{0s}}{T_s}(1-\Gamma_{0s})\phi_{\bm k}
=\sum_sq_s\int J_{0s}g_{s,\bm k}\dd^3v,
\]
\[
\Gamma_{0s}=I_0(b_s)\ee^{-b_s},\qquad
b_s=k_\perp^2T_s/(m_s\Omega_s^2).
\]
The two forms use different distribution variables, so the polarization term appears in different places.
The zero-wave-number potential gauge and particle-number compatibility require separate specification.

Electromagnetic models also determine $A_\parallel$ through parallel current
and include $\delta B_\parallel$, polarization, and magnetization according to their ordering.
The Amp\`ere current must use a push-back consistent with the gyrocenter coordinates, distribution variable, and field equations.
These relations couple the retained velocity distributions to electromagnetic potentials; truncating higher moments is a separate closure choice.

\section{Collisions, Turbulent Transport, and Discretization Levels}
Gyrokinetic collision models must preserve the appropriate particle, momentum, and energy moments
and be compatible with the reduced measure and entropy dissipation. Turbulent particle and heat fluxes arise from correlations between distribution perturbations and radial drifts;
a particle flux has the structure $\langle\int\delta f_s\,v_{E,r}\dd^3v\rangle$.
Specific fluxes depend on the push-back definition and averaging region.

The five-dimensional state still requires discretization in position, parallel velocity, magnetic moment, and time.
Domain extent, velocity-tail truncation, and collision resolution control different errors;
spectral, DG, or particle representations are subsequent numerical choices.

\aiexercise{For a Maxwellian equilibrium and uniform field, calculate $\langle\exp(\ii\bm k_\perp\cdot\bm\rho)\rangle$ and verify $\int F_{0s}J_{0s}^2\dd^3v=n_{0s}\Gamma_{0s}$. Prove that the $h_s$ and $g_s$ quasineutrality forms are equivalent.}
{Starting from~\eqref{eq:slab-gk}, fix periodic boundaries and velocity truncation, and list the data needed to check quasineutrality residuals, free energy, and velocity-space convergence. For further reading, choose one electromagnetic model and map its distribution variable, Jacobian, collision term, and field equations consistently.}
